\begin{textsample}{POS Dim 6 – human – Score 3.00 – mg/mg\_ef\_linguagens\_lp\_13\_p1.txt}  \label{ex:f6_pos_019}
Trata -se, neste campo, de ampliar e qualificar a participação dos \textbf{jovens} nas práticas relativas ao debate de ideias e à atuação política e social, por meio do(a) : Compreensão dos interesses que movem a esfera política em seus diferentes níveis e instâncias, das formas e canais de participação institucionalizados, incluindo os digitais, e das formas de participação não institucionalizadas, incluindo aqui manifestações artísticas e intervenções urbanas. \textbf{Reconhecimento} da importância de se envolver com questões de interesse público e \textbf{coletivo} e compreensão do contexto de promulgação dos \textbf{direitos} humanos, das políticas afirmativas, e das leis de uma forma geral em um estado democrático, como forma de propiciar a vivência democrática em várias instâncias e uma atuação pautada pela ética da responsabilidade ( o outro tem tanto \textbf{direito} a uma vida digna quanto eu tenho ).
Desenvolvimento de habilidades e aprendizagem de procedimentos envolvidos na leitura/escuta e produção de textos pertencentes a gêneros relacionados à discussão e implementação de propostas, à defesa de \textbf{direitos} e a projetos culturais e de interesse público de diferentes naturezas.

% matched lemmas: coletivo, direito, jovem, reconhecimento
\end{textsample}
